\documentclass[12pt,a4paper]{ctexart}
\usepackage{geometry}\geometry{a4paper,margin=2cm,landscape=true}
\usepackage{float,booktabs,longtable,enumitem,hyperref,fancyhdr,xcolor}
\pagestyle{fancy}
\fancyhead[L]{dgiot\_lite — 设备与测点完整目录}
\fancyhead[R]{2026-07-12}
\fancyfoot[C]{\thepage}

\begin{document}

\begin{center}
{\LARGE \textbf{大庆采油厂 CY1C8K 设备与测点完整目录}}\\[6pt]
{\large Oracle 11g @ 11.66.12.129:1521/orcl | DQYTPROD}\\[4pt]
{\normalsize 采集方式: IoMonitor (CommitRealSpan=300ms) | 总览: 966口井 $\cdot$ 44,977测点 $\cdot$ 4,815,777功图}
\end{center}

\section{测点类型编码说明}

\begin{longtable}{p{2.5cm}p{4cm}p{8cm}}
\toprule
\textbf{编码} & \textbf{中文名} & \textbf{说明} \\
\midrule
TGP & 套压 & Tubing Pressure — 油管压力 \\
GYS & 工况参数 & Working Condition Status \\
ZWG & 总无功功率 & Total Reactive Power \\
ZYGX & 总有功功率 & Total Active Power \\
ZHL & 总回流 & Total Backflow \\
CZT & 齿轮状态 & Gear Status \\
ADL/BDL/CDL & A/B/C相电流 & Phase Current \\
ADY/BDY/CDY & A/B/C相电压 & Phase Voltage \\
YIS & 油压仪状态 & Oil Pressure Instrument Status \\
DWL & 低回流 & Low Backflow \\
UWL & 高回流 & High Backflow \\
DCV & 低控制阀 & Low Control Valve \\
UCV & 高控制阀 & High Control Valve \\
CHC & 冲程 & Stroke \\
SLV & 位置 & Position \\
CPV & 控制阀位置 & Control Valve Position \\
RCV & 远程控制阀 & Remote Control Valve \\
TX/YX & 遥信/遥测 & Telemetry Signal \\
FDD & 功图 & Dynamometer Diagram \\
PDL & 泵电流 & Pump Current Load \\
PDY & 泵电压 & Pump Voltage \\
\bottomrule
\end{longtable}

\section{按计量站分布的测点数}

\begin{table}[H]\centering
\begin{tabular}{p{2cm}p{2cm}p{2cm}p{2cm}p{2cm}p{2cm}}
\toprule
\textbf{站号} & \textbf{测点数} & \textbf{站号} & \textbf{测点数} & \textbf{站号} & \textbf{测点数} \\
\midrule
站1 (ATU) & 195 & 站2 & 195 & 站3 & 185 \\
站4 & 180 & 站5 & 178 & 站6 & 172 \\
站7 & 170 & 站8 & 171 & 站9 & 169 \\
站10 & 170 & 站11 & 170 & 站12 & 170 \\
\bottomrule
\end{tabular}
\end{table}

\section{测点类型分布 (4567点)}

\begin{longtable}{p{3cm}p{3cm}p{3cm}p{3cm}p{3cm}}
\toprule
\textbf{类型} & \textbf{数量} & \textbf{类型} & \textbf{数量} & \textbf{类型} & \textbf{数量} \\
\midrule
ATUS (站1全量) & 195 & TX/YX (遥信遥测) & 各187 & RCV (远程阀) & 190 \\
TGP (套压) & 190 & GYS (工况) & 190 & ZWG (总无功) & 190 \\
ZYG (总有功) & 190 & ZHL (总回流) & 190 & CZT (齿轮) & 190 \\
ADL (A相电流) & 190 & BDL (B相电流) & 190 & CDL (C相电流) & 190 \\
ADY (A相电压) & 190 & BDY (B相电压) & 190 & CDY (C相电压) & 190 \\
YIS (油压仪) & 185 & DWL (低回流) & 184 & UWL (高回流) & 184 \\
DCV (低控阀) & 184 & UCV (高控阀) & 184 & CHC (冲程) & 184 \\
SLV (位置) & 184 & CPV (控阀位) & 190 & RCV3/5 & 27 \\
FDD (功图) & 14 & PDL (泵电流) & 21 & PDY (泵电压) & 21 \\
\bottomrule
\end{longtable}

\section{OPC DA 数据 (SYS\_POINTRELATION\_STATION)}

\textbf{OPC DA 数据确实在 Oracle 中！} 之前误判为"没有 OPC 数据"是因为只检查了 SYS\_POINTRELATION\_WELL 表。实际 OPC DA 数据存储在 SYS\_POINTRELATION\_STATION 表中。

\begin{table}[H]\centering
\caption{数据按协议来源分布}
\begin{tabular}{p{4cm}p{3cm}p{3cm}p{4cm}}
\toprule
\textbf{表名} & \textbf{行数} & \textbf{协议来源} & \textbf{说明} \\
\midrule
SYS\_POINTRELATION\_WELL & 4,567 & Modbus TCP & CommBridge 采集 (B1V/B2V/B3V) \\
SYS\_POINTRELATION\_STATION & 40,410 & OPC DA + 注水站 & OPC\_FC\_Client 采集 (DX/JB/Z1/Z2) \\
\midrule
\textbf{合计} & \textbf{44,977} & & \\
\bottomrule
\end{tabular}
\end{table}

\begin{table}[H]\centering
\caption{OPC DA 测点分布 (SYS\_POINTRELATION\_STATION)}
\begin{tabular}{p{4cm}p{3cm}p{4cm}p{3cm}}
\toprule
\textbf{井号前缀} & \textbf{测点数} & \textbf{示例路径} & \textbf{来源} \\
\midrule
DX (DX8ZRZ等) & 16,375 & /CY1C8K/DX8ZRZ/ZD010825... & OPC DA \\
JB (JB1V2等) & 969 & /CY1C8K/JB1V2/ZD010826... & OPC DA \\
Z2 (注水站) & 12,612 & /CY1C8K/SY217Z2VJLJV1/... & 注水站 \\
Z1 (注水站) & 6,202 & /CY1C8K/SY217Z1VJLJV1/... & 注水站 \\
其他 & 4,252 & /CY1C8K/SY... & 混合 \\
\midrule
\textbf{OPC DA 合计} & \textbf{17,344} & & \\
\bottomrule
\end{tabular}
\end{table}

\section{Modbus TCP 数据 (SYS\_POINTRELATION\_WELL)}

全部 4,567 条测点来自 CommBridge.exe (Modbus TCP)，井号前缀为 B1V/B2V/B3V。

\section{最新运行数据 (2026-07-12 05:10)}

\begin{longtable}{p{2.5cm}p{3cm}p{3cm}p{3cm}p{4cm}}
\toprule
\textbf{井号} & \textbf{数据时间} & \textbf{运行率} & \textbf{今日运行} & \textbf{累计运行} \\
\midrule
8035 & 2026/7/12 5:10 & 21.53\% & 310min & — \\
8127 & 2026/7/12 5:10 & 21.53\% & 310min & — \\
8028 & 2026/7/12 5:10 & 21.53\% & 310min & — \\
8128 & 2026/7/12 5:10 & 21.53\% & 310min & — \\
8171 & 2026/7/12 5:10 & 21.53\% & 310min & — \\
8530 & 2026/7/12 5:10 & 21.53\% & 310min & — \\
8012 & 2026/7/12 5:10 & 21.53\% & 310min & — \\
8014 & 2026/7/12 5:10 & 18.06\% & 260min & — \\
8175 & 2026/7/12 5:10 & 0\% & — & — \\
8011 & 2026/7/12 5:10 & 20.14\% & 290min & — \\
\bottomrule
\end{longtable}

\section{井号清单 (966口)}

\noindent 完整井号清单见配套CSV文件: \texttt{docs/wells.csv} (966行)

\vspace{8pt}
\noindent 抽样展示 (前100口):

\begin{longtable}{p{3cm}p{2.5cm}p{3cm}p{2.5cm}p{3cm}p{2.5cm}}
\toprule
\textbf{井号} & \textbf{频率} & \textbf{井号} & \textbf{频率} & \textbf{井号} & \textbf{频率} \\
\midrule
保5-23向1 & 1.70 & 保5-31 & — & 保5-23向4 & 3.10 \\
保5-23向3 & 3.80 & 保5-23向2 & 3.20 & 保7-02向3 & 3.20 \\
保7-02向2 & 3.00 & 保7-02向1 & 2.00 & 保7-02 & 2.60 \\
保5-31向4 & — & 葡2-27向2 & — & 葡5-23向1 & — \\
\bottomrule
\end{longtable}

\vspace{12pt}
\begin{center}
\textit{完整数据请参见: \texttt{docs/wells.csv} (966行) 和 \texttt{docs/points.csv} (4567行)}
\end{center}

\section{数据源说明}

\begin{itemize}[leftmargin=*]
  \item \textbf{数据库}: Oracle 11g @ 11.66.12.129:1521/orcl
  \item \textbf{用户}: DQYTPROD
  \item \textbf{采集软件}: 力控 ForceControl 7.x / IoMonitor v6.0.0.1
  \item \textbf{采集方式}: CommBridge.exe (Modbus TCP) → IoMonitor → Oracle
  \item \textbf{提交频率}: 300ms (CommitRealSpan)
  \item \textbf{场站}: CY1C8K (采油一厂八矿) — 4372 测点
  \item \textbf{协议}: Modbus TCP (Port 53001) — 连接 80+ 井口 RTU (11.248.x.x / 11.249.x.x)
  \item \textbf{导出时间}: 2026-07-12 05:12 CST
  \item \textbf{工具}: oracle\_reader.py (独立 Python 脚本, WinRM 中继)
\end{itemize}

\end{document}
